\sgaStableTarget{sga6.e05.para.3-10}3.10
Let @@M0@@ be an augmented @@M1@@-@@M2@@-algebra. We define a ring filtration on @@M3@@ as follows: for every @@M4@@, we denote by @@M5@@ the ideal generated by the products:
\sgaStableTarget{sga6.e05.eq.3-10-1}@@M6@@
@@M7@@ being the augmentation).

Suppose that @@M8@@ is generated as an abelian group by a family of elements @@M9@@; then @@M10@@ is generated as a @@M11@@-module by the products:
\sgaStableTarget{sga6.e05.eq.3-10-2}@@M12@@
Indeed, if @@M13@@ is an element of @@M14@@ with zero augmentation and we write @@M15@@, we have @@M16@@, hence @@M17@@; furthermore, if @@M18@@, we can write @@M19@@; using relations (3.4.1), we therefore see that the elements of @@M20@@ are linear combinations with coefficients in @@M21@@ of the elements (3.10.2); since these belong to @@M22@@, we obtain the required result.

The filtration thus defined is a ring filtration; it is decreasing, and we have:
@@M23@@
Moreover, it is functorial in @@M24@@: every @@M25@@-homomorphism of augmented @@M26@@-@@M27@@-algebras is compatible with the filtrations. When we regard an augmented @@M28@@-ring as endowed with a filtration, this will mean the filtration above, unless expressly stated otherwise.

